% Gerado a partir de MONOGRAFIA_PASE.md. O texto foi transposto sem alterações.

\chapter{Trabalhos Relacionados}\label{cap:3}

\section{Introdução}\label{sec:3.1}

Este capítulo apresenta os trabalhos que mais se aproximam da proposta. A seleção partiu das quatro frentes que compõem a Arquitetura PASE. As duas primeiras são (i) a supervisão humana no fluxo de desenvolvimento com agentes e (ii) a automação do ciclo de vida por sistemas multiagente. As duas seguintes são (iii) a originação de demandas a partir de fontes não estruturadas e (iv) a personalização de interfaces em tempo de execução. Para cada frente, foi escolhido o trabalho mais próximo da PASE entre os publicados em veículos revisados por pares de computação, com prioridade para propostas que descrevem um sistema implementado e avaliado, e não apenas uma visão. Para cada trabalho, descrevem-se sua proposta, os pontos de convergência com a PASE e as diferenças que delimitam a contribuição deste estudo. A Seção~\ref{sec:3.6} consolida a comparação no Quadro~\ref{quadro:2}.

\section{Supervisão Humana: HULA}\label{sec:3.2}

O trabalho mais próximo da PASE em termos de supervisão humana é o HULA (\textit{Human-in-the-Loop LLM-based Agents}), proposto por \citeonline{takerngsaksiri2025hula} e implantado em produção no Jira, da Atlassian. O \textit{framework} organiza-se em três agentes que cooperam para resolver uma \textit{issue}. O \textit{AI Planner Agent} localiza os arquivos afetados e gera um plano de codificação, e o \textit{AI Coding Agent} implementa o plano. O \textit{Human Agent} é o engenheiro que revisa, refina e aprova o plano e o código antes da abertura de um \textit{pull request}. Em dois meses de uso, os profissionais acionaram o HULA em 663 \textit{issues}; foram gerados 527 planos, 433 deles aprovados, e 56 \textit{pull requests} acabaram integrados.

Essa sequência corresponde ao trecho da PASE que vai da Camada de Orquestração à Camada de Execução, em especial ao Documento de Planejamento Técnico, exigido antes da escrita de código, e à homologação humana com autoridade final. Os próprios resultados do HULA reforçam a necessidade dessa supervisão: na pesquisa com 109 engenheiros, apenas 33\% concordaram que o código gerado resolvia a \textit{issue} sem intervenção humana, e 54\% discordaram de que ele estivesse livre de defeitos \cite{takerngsaksiri2025hula}. O HULA, contudo, parte de uma \textit{issue} já formulada e priorizada por pessoas. Ele não trata a originação nem a filtragem de demandas, não modela interdependências entre tarefas, não guarda memória das rejeições e não atua sobre o software em tempo de execução. A PASE pode ser entendida como uma generalização do princípio do HULA, que desloca a supervisão humana de pontos isolados de revisão para um plano de controle que acompanha a tarefa da captação à personalização em produção.

\section{Automação do Ciclo de Vida: MetaGPT}\label{sec:3.3}

Na amplitude do ciclo de vida coberto, destaca-se o MetaGPT \cite{hong2024metagpt}, o \textit{framework} multiagente que mais se aproxima da PASE. Nele, agentes assumem papéis de uma equipe de software (gerente de produto, arquiteto, gerente de projeto, engenheiro e engenheiro de qualidade) e seguem Procedimentos Operacionais Padronizados (\textit{Standard Operating Procedures}, SOPs). Um requisito de uma linha é progressivamente transformado em documento de requisitos, projeto de sistema, lista de tarefas, código e testes, com artefatos estruturados trocados entre os papéis, em linha de montagem, para reduzir a propagação de alucinações. Os autores relatam 85,9\% e 87,7\% de acerto na primeira tentativa nos \textit{benchmarks} HumanEval e MBPP, e um custo de revisão humana menor que o de abordagens anteriores baseadas em conversa entre agentes.

A decomposição de uma demanda em etapas com artefatos intermediários verificáveis e a separação entre quem planeja e quem executa são princípios compartilhados pela PASE, cuja Camada de Execução também produz um documento de planejamento antes da implementação. A diferença central está no papel da pessoa e na origem das demandas. O MetaGPT busca simular uma organização de software inteira com agentes e concentra a intervenção humana no requisito inicial. A PASE, por sua vez, mantém o gestor como responsável pela priorização, pela aprovação do plano e pela homologação, e trata a estratégia de coordenação entre agentes como detalhe interno do Executor. Essa escolha está em consonância com a evidência de que a eficácia de sistemas multiagente depende da adequação entre a estrutura de coordenação e a tarefa, e não do número de papéis envolvidos \cite{kim2026scaling}. Além disso, o MetaGPT não prevê isolamento obrigatório da execução, controle de custo, auditoria recorrente da base de código nem adaptação em tempo de execução.

\section{Originação de Demandas: LLM-Cure}\label{sec:3.4}

Quanto à originação e à qualificação de demandas, a proposta mais próxima da Camada de Tarefas da PASE é o LLM-Cure, de \citeonline{assi2026llmcure}. A abordagem usa modelos de linguagem para extrair as funcionalidades mencionadas em avaliações de usuários em lojas de aplicativos. Diante de uma reclamação sobre determinada funcionalidade, recupera avaliações bem pontuadas de aplicativos concorrentes sobre a mesma funcionalidade e sintetiza sugestões concretas de melhoria. A avaliação, feita sobre 1.056.739 avaliações de setenta aplicativos Android, indicou que, em média, 73\% das sugestões geradas foram posteriormente implementadas pelos próprios desenvolvedores, o que evidencia a viabilidade de converter \textit{feedback} informal e dados de concorrentes em propostas acionáveis.

Esse desenho corresponde, na PASE, à captação passiva do fluxo de agentes externos e à prospecção competitiva do fluxo da inteligência artificial. O LLM-Cure, entretanto, restringe-se a uma única fonte de insumos e encerra seu ciclo na recomendação ao desenvolvedor. Ele não contempla a normalização de entradas heterogêneas, a deduplicação contra o histórico de funcionalidades e de propostas descartadas, a verificação de coerência com as regras de negócio nem a atribuição de metadados de impacto, complexidade e confiança. A PASE estende essa linha ao tratar a mineração de \textit{feedback} e a comparação com concorrentes como apenas duas das fontes de um funil unificado de filtragem progressiva, cujo produto é um \textit{backlog} qualificado e pronto para priorização sob supervisão humana, e não uma lista de sugestões isoladas.

\section{Personalização em Tempo de Execução: Interfaces Generativas e Maleáveis}\label{sec:3.5}

A personalização em tempo de execução encontra seu precedente acadêmico mais direto em \citeonline{cao2025malleable}, apresentado na CHI 2025. Os autores propõem interfaces generativas e maleáveis. Em vez de telas estáticas definidas em tempo de projeto, a interface é derivada de um modelo de dados orientado à tarefa, gerado por um LLM a partir do \textit{prompt} do usuário e composto pelas entidades, relações e propriedades relevantes à atividade em curso. Esse modelo evolui à medida que o usuário refina sua intenção, e cada alteração se propaga para a interface e para os dados subjacentes. O protótipo Jelly e o estudo com usuários indicaram que os participantes valorizaram a combinação entre elementos estáticos previsíveis e elementos gerados sob demanda.

Esse resultado converge com a decisão da PASE de preservar o comportamento nativo e abrir à personalização apenas áreas delimitadas. A separação entre uma representação estruturada do estado, que o modelo gera, e sua renderização, que o sistema controla, é também o fundamento dos \textit{payloads} declarativos das Zonas Dinâmicas. O trabalho de Cao, Jiang e Xia, porém, parte sempre de uma intenção explícita do usuário e corresponde apenas ao que a PASE denomina Zonas Dinâmicas Sob Demanda.

Não há, no trabalho de Cao, Jiang e Xia, equivalente às Zonas Dinâmicas Autônomas, acionadas pelo UPA a partir de telemetria, nem ao Avaliador de Zonas Dinâmicas, que decide, com base no padrão das demandas recebidas, quais módulos devem se tornar personalizáveis. Tampouco são tratados os mecanismos de governança próprios de um software comercial multiusuário, como contratos rígidos por zona com retorno ao padrão nativo, acesso a dados sob menor privilégio e filtragem prévia de injeção de \textit{prompt}. A PASE incorpora a maleabilidade proposta pelos autores, mas a subordina a uma camada de governança que define onde, quando e sob quais limites a geração pode ocorrer.

\section{Considerações Finais}\label{sec:3.6}

Os trabalhos analisados mostram que cada frente da PASE possui precedentes, mas que nenhum deles a cobre por completo, como resume o Quadro~\ref{quadro:2}. O HULA insere a pessoa na revisão de plano e de código, mas não na priorização nem na memória das decisões rejeitadas. O MetaGPT decompõe o ciclo de vida em papéis e artefatos, mas sem isolamento de execução, controle de custo ou auditoria da base de código. O LLM-Cure converte \textit{feedback} disperso em sugestões, mas sem normalização de múltiplas fontes, deduplicação contra o histórico ou qualificação por metadados. E o trabalho de Cao, Jiang e Xia torna interfaces maleáveis pela intenção do usuário, mas sem gatilho por telemetria, sem critério para decidir quais áreas devem ser personalizáveis e sem contratos que limitem a geração.

\begin{quadro}[htb]
\centering
\caption{Comparação entre os trabalhos relacionados e a Arquitetura PASE}
\label{quadro:2}
\footnotesize
\renewcommand{\arraystretch}{1.25}
\begin{tabularx}{\textwidth}{|X|X|X|X|X|X|}
\hline
\textbf{Capacidade} & \textbf{HULA} & \textbf{MetaGPT} & \textbf{LLM-Cure} & \textbf{Cao, Jiang e Xia} & \textbf{PASE} \\
\hline
Originação de demandas de múltiplas fontes & Não & Não & Parcial (uma fonte) & Não & Sim \\
\hline
Filtragem com memória de implementados e descartados & Não & Não & Não & Não & Sim \\
\hline
Decisão humana sobre o plano e sobre a entrega & Sim & Parcial (requisito inicial) & Não se aplica & Não se aplica & Sim \\
\hline
Dependências entre tarefas e controle de custo & Não & Não & Não & Não & Sim \\
\hline
Execução isolada e verificação por código & Parcial & Parcial & Não se aplica & Não & Sim \\
\hline
Personalização em tempo de execução governada por contrato & Não & Não & Não & Parcial (sem contrato) & Sim \\
\hline
\end{tabularx}
\fonte{elaborado pelos autores (2026).}
\end{quadro}

A contribuição da PASE, portanto, não se resume a integrar essas frentes. Dentro de cada camada, ela introduz mecanismos que os trabalhos relacionados não contemplam e, entre as camadas, subordina todos eles a um mesmo regime de autonomia supervisionada, rastreabilidade e contenção da generatividade.
